\subsection{Agreement-based cascading evaluation}\label{subsec:eval-abc}

The agreement-based cascading evaluation compares the calibrated cascading system to the single-model baselines. The cascade configuration is selected on the calibration set of Subsection~\ref{subsec:knowledge-extraction-calibration-set}; its strict-F1 and cost are reported on that calibration cluster, and the calibrated cascade is then evaluated on the held-out evaluation set (Subsection~\ref{subsec:knowledge-extraction-eval-set}) and the results are reported separately as \textbf{RQ5} (Subsection~\ref{subsec:eval-heldout-generalization}).

The cascade configuration is selected using the offline replay procedure described in Subsection~\ref{subsec:eval-reprod-offline-replay}, in a screen-then-refine sequence: a screen over every combination of embedder (gemini, openai), agreement scorer (embedding, hybrid), and alignment strategy (greedy, Hungarian, soft\_max); a weight refinement of surviving structures (the hybrid weights and the soft\_max parameters); and a later experimental pass over the agreement threshold $\theta \in \{0.3, 0.4, \dots, 0.9\}$ and the rejection threshold $\theta_{\mathrm{reject}} \in \{0.0, 0.1, 0.2\}$. The safety gate policies of Subsection~\ref{subsec:abc-safety-gates} are then evaluated as a separate $2\times2$ ablation (single-voter policy $\times$ polarity-conflict forcing) at the pinned configuration, after the joint sweep on the other variables is complete. The best configuration values found on the calibration set are reported as empirical results in Chapter~\ref{chapter:Results}. 


% which sweeps the agreement threshold $\theta \in \{0.3, 0.4, \dots, 0.9\}$, the rejection threshold $\theta_{\mathrm{reject}} \in \{0.0, 0.1, 0.2\}$, and the different agreement scoring strategies: embedding and hybrid, described in Subsection~\ref{subsec:abc-agreement-scoring}. 

The cascade with its calibrated configuration values is compared to a set of single-model baselines. All comparisons are done under identical matcher settings on the calibration-cluster's silver findings, so that the comparisons differ only in the generated pipeline outputs. Strict F1-score of each system, as well as cost (per-chunk escalation rate and incurred cost), is reported in Chapter~\ref{chapter:Results}.

\pinktodo{Add leave-one-voter-out ablations.}

\pinktodo{Add bootstrap experiments to the code.}

\pinktodo{When the experiments are ready to run, explain what they are in here as a final paragraph.}